% f(x) = x^3 - 6x^2 + 9x + 1: maximum (1,5), minimum (3,1).  Flat tangents drawn at both;
% gradient signs from f'(x) = 3(x-1)(x-3): f'(0)=9>0, f'(2)=-3<0, f'(4)=9>0.  Caller sets \figunit.
\begin{tikzpicture}[x=\figunit, y=0.55\figunit, line cap=round]
  \draw[-{Stealth[length=2mm]}, line width=0.6pt] (-0.6,0) -- (5.6,0) node[right, inner sep=1pt] {$x$};
  \draw[-{Stealth[length=2mm]}, line width=0.6pt] (0,-0.5) -- (0,7.3) node[above, inner sep=1pt] {$y$};
  \begin{scope}
    \clip (-0.6,-0.5) rectangle (4.8,7.2);
    \curve{-0.25}{4.45}{\x*\x*\x-6*\x*\x+9*\x+1}
  \end{scope}
  \draw[line width=0.8pt, dashed] (0.4,5) -- (1.6,5);
  \draw[line width=0.8pt, dashed] (2.4,1) -- (3.6,1);
  \fill (1,5) circle (2.2pt); \fill (3,1) circle (2.2pt);
  \node[above, inner sep=2pt] at (1,5.1) {maximum};
  % 'minimum' right of the flat tangent: at x = 3.7 the curve is at 2.8, above the label
  \node[anchor=west, inner sep=2pt] at (3.65,1) {minimum};
  % gradient signs, placed beside the curve
  \node at (0.6,2.2) {\textbf{$+$}};   % below the rising curve (f(0.6) = 4.5), clear of the y-axis
  \node at (2.35,4.3) {\textbf{$-$}};
  \node at (4.35,3.1) {\textbf{$+$}};
\end{tikzpicture}
